\documentclass[10pt,a4paper]{article}
\usepackage[T1]{fontenc}
\usepackage[utf8]{inputenc}
\usepackage{graphicx}
\usepackage{polski}
\usepackage[left=1in,right=1in,top=1in,bottom=1in]{geometry}

\bibliographystyle{unsrt}

\begin{document}
\begin{titlepage}
	\centering
	\includegraphics[width=10cm]{media/PB_logo}\par\vspace{1cm}
	{\LARGE\bfseries Sztuczna inteligencja \par}
	\vspace{0.5cm}
	
	{\LARGE\bfseries Sprawozdanie z projektu końcowego\par}
	\vspace{1.5cm}
	
	{\Large Temat: Prezentacja możliwości biblioteki scikit-learn w projektach informatycznych, wymagających użycia drzew decyzyjnych.\par}
	\vspace{3cm}
	
	\raggedright
	
	
 	{\large Wykonujący projekt: \textbf{Adam Bajguz}\\
 			\makebox[0pt][l]{}\phantom{Wykonujący projekt: }\textbf{Magdalena Kalisz}\par}
 	\vspace{1cm}
 	
 	{\large Studia dzienne \par
	Kierunek: Informatyka \par
	Semestr: IV \hspace{2.5cm} Grupa zajęciowa: \textbf{PS1} \par
	Prowadzący ćwiczenie: \textbf{mgr inż. Dariusz Jankowski} \par}
 	
	\vfill
		
	\begin{tabular}{c}
	\today\\
	Data wykonania projektu
	\end{tabular}
	
	\begin{flushright}
	\begin{tabular}{c}
	..............................................\\
	Data i podpis prowadzącego
	\end{tabular}
	\end{flushright}

\end{titlepage}

\tableofcontents
\newpage

\section{Wstęp}
rtt r  tt\cite{test} t t t
\section{Idea drzew decyzji}
\subsection{Drzewa klasyfikacyjne i regresyjne}
\subsection{Drzewa decyzyjne w teorii decyzji}
\subsection{Drzewa decyzyjne w typowych zastosowaniach CI: klasyfikacja danych i maszynowym uczeniu się}
\subsection{Budowa drzew decyzyjnych}
\subsection{Cechy drzew decyzyjnych}
\subsection{Metody poprawy wyników uzyskiwanych dzięki drzewom decyzyjnym}
\subsection{Przycinanie drzew decyzyjne}
\subsection{Inne metody poprawy wyników}

\section{Algorytmy budowania drzew decyzyjnych}
\subsection{Algorytm ID3}
\subsection{Algorytm C4.5}
\subsection{Algorytm C5.0}
\subsection{Algorytm CART}

\section{Entropia}

\section{Funkcje podziału}
\subsection{GINI}
\subsection{Information gain}
\subsection{Gain ratio}
\subsection{Random}

\section{Biblioteka scikit-learn}

\section{Opis programu}

\section{Wnioski}

\section{Omówienie literatury i bieżącego stanu wiedzy w wybranej dziedzinie SI}

\section{Omówienie wykorzystanej biblioteki i ewentualnych modyfikacji kodu źródłowego}

\section{Opisanie danych wykorzystanych do testowania działania metody wybranej biblioteki}

\section{Omówienie głównych części programu/skryptu i wyników}

\section{Wnioski końcowe odnośnie biblioteki i projektu}

\addcontentsline{toc}{section}{Literatura}
\bibliography{Projekt}

\end{document}